% =============================================================================
% preguntas/p5-kitty-problem/solucion.tex — P5 (issue #9)
% Integrante: Ernesto A. Soler Choong
% Fragmento puro: se incluye con \include desde SemArq1_TeamBales.tex.
% Prohibido: \documentclass, \begin{document}, \usepackage, preámbulo.
% Fuente de trabajo: solucion.md (misma carpeta).
% Citas: biblatex (APA) con páginas impresas; Richards & Ford (2020), Cockburn
% (2005) y Palermo (2008) no tienen paginación y se citan por capítulo/sección.
% Claves de core.bib usadas: martin2017, richards2015, richardsford2020,
% delatorre2010, cockburn2005. palermo2008 está declarada en
% preguntas/p3-ddd-bounded-context/referencias.bib (o en core.bib, si se mueve).
% Paquetes que debe cargar el preámbulo común: xcolor, tikz, listings y
% enumitem (con babel spanish + es-nolists).
% =============================================================================

% --- Definiciones auxiliares. Son globales (siguen activas en los capítulos
% posteriores); son idénticas a las de P3, así que no entran en conflicto. ---
\providecolor{azul}{HTML}{1F4E79}
\providecolor{naranja}{HTML}{C55A11}
\providecolor{gris}{HTML}{595959}
\usetikzlibrary{arrows.meta}
\lstdefinestyle{codigo}{
  basicstyle=\ttfamily\footnotesize,
  keywordstyle=\color{azul}\bfseries,
  commentstyle=\color{gris}\itshape,
  stringstyle=\color{naranja},
  frame=lines, framesep=4pt, columns=fullflexible, keepspaces=true,
  showstringspaces=false, breaklines=true, xleftmargin=1em,
  literate={á}{{\'a}}1 {é}{{\'e}}1 {í}{{\'i}}1 {ó}{{\'o}}1 {ú}{{\'u}}1
           {ñ}{{\~n}}1 {Á}{{\'A}}1 {É}{{\'E}}1 {Í}{{\'I}}1 {Ó}{{\'O}}1
           {Ú}{{\'U}}1 {Ñ}{{\~N}}1
}
\renewcommand{\lstlistingname}{Listado}

\chapter{P5: Análisis de un Problema Real: El Kitty Problem y el Principio de Abierto/Cerrado}
\label{sec:p5}

\section{Planteamiento}
\label{sec:p5-planteamiento}

El enunciado presenta el \emph{Kitty Problem}: un requisito transversal, la entrega de
gatitos, afecta a varios servicios ya existentes de un sistema de taxis (\texttt{TaxiUI},
\texttt{TaxiFinder}, \texttt{TaxiSelector}, etc.) y obliga a cambiarlos todos de forma
coordinada\footnote{El enunciado sitúa el ejemplo en \textit{Fundamentals of Software
Architecture} \parencite{richardsford2020}. En la bibliografía consultada, el caso aparece
en \textcite[cap.~27]{martin2017}, que es la fuente que se cita en esta respuesta.}. La
pregunta pide:

\begin{enumerate}[label=(\alph*)]
  \item describir el problema y explicar por qué dividir el sistema en servicios basados
        en verbos o entidades funcionales (como \texttt{TaxiFinder} o
        \texttt{TaxiSelector}) llevó a él;
  \item explicar cómo, con los principios de DDD y de la Arquitectura Onion/Clean, podría
        haberse diseñado el sistema para que un cambio así afectara a un solo componente
        o servicio, y
  \item proponer una solución en la que componentes como \texttt{Kittens} y
        \texttt{Rides} se añadan como \emph{plugins}, explicando cómo respeta el Principio
        de Abierto/Cerrado (OCP) y qué papel desempeñan los adaptadores.
\end{enumerate}

\section{El Kitty Problem}
\label{sec:p5-problema}

Martin parte de un agregador de taxis, ya presentado en el capítulo sobre el Principio de
Sustitución de Liskov, que conoce a muchos proveedores de taxis de una ciudad y permite a
los clientes pedir viajes \parencite[80, 242]{martin2017}. Para que fuera escalable, el
sistema se construyó con muchos microservicios, y el personal se dividió en equipos
pequeños, cada uno responsable de unos pocos servicios \parencite[242]{martin2017}:
\begin{itemize}
  \item \texttt{TaxiUI} atiende a los clientes, que piden taxis desde el móvil;
  \item \texttt{TaxiFinder} consulta los inventarios de los proveedores y determina los
        taxis candidatos;
  \item \texttt{TaxiSelector} elige uno según los criterios del cliente (coste, tiempo,
        lujo);
  \item \texttt{TaxiDispatcher} encarga el taxi elegido.
\end{itemize}

Tras más de un año en producción, el departamento de \textit{marketing} anuncia un servicio
de entrega de gatitos. Habrá puntos de recogida en la ciudad; un taxi cercano recogerá el
gatito y lo llevará a la dirección indicada. Por ahora, un proveedor ha aceptado
participar; es probable que otros se sumen y que otros lo rechacen. Los conductores
alérgicos a los gatos nunca deben ser elegidos, y un vehículo que haya transportado
gatitos en los últimos tres días no debe asignarse a clientes que declaren esa alergia
\parencite[243]{martin2017}. La consecuencia es que \emph{todos} los servicios tienen que
cambiar y su desarrollo y despliegue deben coordinarse: los servicios están acoplados y
no pueden desarrollarse, desplegarse ni mantenerse de forma independiente
\parencite[243]{martin2017}. La Figura~\ref{fig:p5-problema} resume la situación.

\begin{figure}[htbp]
  \centering
  \resizebox{\ifdim\width>\linewidth\linewidth\else\width\fi}{!}{%
  \begin{tikzpicture}[
      srv/.style={draw=azul, fill=azul!6, rounded corners=3pt, thick, font=\small\ttfamily,
                  minimum width=3.0cm, minimum height=1.1cm, align=center},
      flujo/.style={-{Stealth[length=2.4mm]}, thick, black!70},
      corte/.style={-{Stealth[length=2.2mm]}, thick, dashed, naranja}
    ]
    \node[srv] (ui)  at (0,0)    {TaxiUI};
    \node[srv] (fin) at (3.9,0)  {TaxiFinder};
    \node[srv] (sel) at (7.8,0)  {TaxiSelector};
    \node[srv] (dis) at (11.7,0) {TaxiDispatcher};
    \draw[flujo] (ui) -- (fin);
    \draw[flujo] (fin) -- (sel);
    \draw[flujo] (sel) -- (dis);
    \node[draw=naranja, fill=naranja!10, rounded corners=3pt, thick, font=\small,
          minimum width=13.6cm, minimum height=0.9cm, align=center] (req) at (5.85,2.1)
      {\textbf{Requisito transversal: entrega de gatitos}\\[1pt]
       \footnotesize puntos de recogida · proveedores participantes · conductores y clientes alérgicos};
    \foreach \n in {ui,fin,sel,dis} {
      \draw[corte] (req.south -| \n) -- node[right, font=\scriptsize, naranja]{cambia} (\n.north);
    }
    \node[font=\scriptsize, gris] at (5.85,-1.0)
      {descomposición funcional: cada servicio es un paso (un verbo) del flujo «pedir un taxi»};
  \end{tikzpicture}%
  }
  \caption{El Kitty Problem: el nuevo requisito atraviesa los cuatro servicios, que deben
           cambiar y desplegarse de forma coordinada. Elaboración propia a partir de
           \textcite[242--243]{martin2017}.}
  \label{fig:p5-problema}
\end{figure}

\section{Por qué la división por verbos o entidades llevó al problema}
\label{sec:p5-causas}

\paragraph{La descomposición funcional es vulnerable a los cambios transversales.}
Cada servicio implementa un paso del flujo: encontrar, seleccionar, despachar. Un
requisito que introduce un \emph{tipo nuevo} de pedido no cae en un paso, sino en todos.
Martin lo resume así: todo sistema se enfrenta a los aspectos transversales, y las
descomposiciones funcionales como la de este ejemplo son muy vulnerables a las
funcionalidades que atraviesan todos esos comportamientos \parencite[243]{martin2017}.

\paragraph{Los servicios no estaban realmente desacoplados.}
Separar en procesos distintos desacopla las variables, pero no los datos: si se añade un
campo a un registro que circula entre servicios, todos los que lo usan deben cambiar y
acordar su interpretación. En esa medida, el desarrollo, el despliegue y la operación
deben coordinarse \parencite[240--241]{martin2017}. En el ejemplo, la lista de candidatos
que \texttt{TaxiFinder} deja para \texttt{TaxiSelector} tendría que incorporar la
información sobre gatitos y alergias.

\paragraph{Las fronteras siguen los pasos del flujo, no el dominio.}
\textcite[cap.~8]{richardsford2020} distinguen la partición técnica, por capas como
presentación, reglas de negocio o persistencia, de la partición por dominio. Con la
primera, cada flujo de negocio queda repartido por todos los componentes; con la segunda,
los componentes de nivel superior se organizan por dominio o flujo, lo que refleja mejor
los cambios que suelen ocurrir. Los pasos del flujo no son capas técnicas en sentido
estricto, pero el efecto es análogo: la nueva capacidad de negocio queda repartida entre
todos ellos. Dividir por entidades tampoco lo resolvería. Con un servicio por entidad
(taxis, pedidos, proveedores), la entrega de gatitos seguiría afectando a todos, y es el
antipatrón \emph{entity trap}, que confunde las relaciones entre datos con los flujos de
trabajo de la aplicación \parencite[cap.~8]{richardsford2020}.

\paragraph{Se violó el Principio de Cierre Común (CCP).}
El CCP pide reunir en un mismo componente lo que cambia por las mismas razones y al mismo
tiempo, de modo que un cambio de requisitos afecte al menor número posible de
componentes \parencite[105--106]{martin2017}. Todas las reglas de los gatitos cambian
juntas, pero el diseño las dispersa por cuatro servicios. Martin vincula expresamente
este principio con el Kitty Problem \parencite[106]{martin2017}.

\section{Diseño según DDD y la Arquitectura Onion/Clean}
\label{sec:p5-diseno}

La clave es trazar las fronteras sobre el \emph{eje de cambio} correcto: los componentes
de cada lado de una frontera cambian a ritmos distintos y por razones distintas
\parencite[173]{martin2017}. En este sistema, lo que varía no son los pasos del flujo,
sino el \emph{tipo de servicio de transporte}: pasajeros, gatitos y los que se añadan en
el futuro. El diseño se construye sobre tres ideas.

\paragraph{1. Un núcleo de políticas estables, protegido por la Regla de Dependencia.}
Las políticas genéricas de encontrar, seleccionar y despachar un vehículo forman el
núcleo, expresado con clases abstractas e interfaces. Las dependencias del código fuente
solo pueden apuntar hacia adentro, hacia las políticas de más alto nivel
\parencite[203]{martin2017}. Es la misma idea de la Arquitectura Onion: las interfaces se
definen en el núcleo y las implementaciones que dependen de detalles viven en los bordes
\parencite{palermo2008}. Si el núcleo debe quedar protegido de los cambios en las
variantes, son las variantes las que deben depender del núcleo
\parencite[73--74]{martin2017}.

\paragraph{2. Contextos por capacidad de negocio (DDD).}
Cada tipo de servicio se convierte en un componente que agrupa todo su modelo:
\texttt{Rides} para los viajes de pasajeros y \texttt{Kittens} para las entregas de
gatitos. Desde DDD, cada uno se comporta como un contexto delimitado: lo relativo a su
dominio es visible dentro y opaco fuera, y cada contexto define sus propias entidades en
lugar de compartir una clase única \parencite[cap.~7]{richardsford2020}. Cada contexto
tiene además su propio Lenguaje Ubicuo, el vocabulario de los expertos en el que se
expresa el modelo \parencite[21]{delatorre2010}:
\begin{itemize}
  \item en \texttt{Rides}, el cliente es un \emph{pasajero} y el pedido es un \emph{viaje};
  \item en \texttt{Kittens}, el cliente es quien recibe el gatito y el pedido es una
        \emph{entrega} con un punto de recogida.
\end{itemize}
Las alergias o el historial de entregas solo existen en el lenguaje de \texttt{Kittens},
así que viven en su modelo y no en el núcleo. La Guía de Arquitectura N-Capas describe
esta organización como módulos verticales por área de negocio, de alta cohesión y bajo
acoplamiento, que se comunican mediante interfaces y cuyos nombres forman parte del
Lenguaje Ubicuo \parencite[71--72]{delatorre2010}. El resultado es una partición por
dominio en la que cada capacidad es funcionalmente cohesiva
\parencite[cap.~8; cap.~17, «Granularity»]{richardsford2020}, y que
cumple el CCP, porque todo lo que cambia junto vive junto \parencite[105]{martin2017}.

\paragraph{3. Las fronteras atraviesan los servicios.}
Los servicios no son lo que define la arquitectura. Para tratar los aspectos
transversales, cada servicio debe tener una arquitectura interna de componentes que siga
la Regla de Dependencia: las fronteras arquitectónicas no caen entre servicios, sino que
los atraviesan \parencite[246--247]{martin2017}. En el despliegue, los servicios pueden
seguir existiendo como antes. Cada uno tiene su diseño interno de componentes, y la nueva
funcionalidad se añade como componentes derivados que se incorporan a la ruta de carga de
cada servicio \parencite[246]{martin2017}. La Figura~\ref{fig:p5-solucion} muestra, por
tanto, una vista lógica: el código de \texttt{Kittens} puede repartirse entre los
servicios que lo necesiten.

Con este diseño, añadir la entrega de gatitos consiste en añadir el componente
\texttt{Kittens} y ajustar la interfaz de usuario y su configuración. El núcleo y
\texttt{Rides} no se modifican. Conviene precisar que, incluso en la solución de Martin,
\texttt{TaxiUI} también cambia \parencite[244]{martin2017}: el cambio no afecta
\emph{literalmente} a un único componente, sino que queda confinado a un componente nuevo
más un ajuste en la capa más externa. Es lo esperable, porque el cierre total es
inalcanzable, por lo que el cierre debe ser estratégico \parencite[106]{martin2017}.

\section{Kittens y Rides como plugins: OCP y adaptadores}
\label{sec:p5-plugins}

\paragraph{La solución de Martin.}
Al aplicar SOLID, Martin conserva gran parte de la lógica de los servicios originales en
clases base abstractas. Lo específico de los viajes se extrae a un componente
\texttt{Rides}, y la nueva capacidad se sitúa en un componente \texttt{Kittens}. Ambos
extienden las clases abstractas mediante \emph{Template Method} o \emph{Strategy},
siguen la Regla de Dependencia y son creados por factorías bajo el control de la interfaz
de usuario. Al implementar los gatitos cambia \texttt{TaxiUI}, pero nada más: se añade un
nuevo \emph{jar}, \emph{Gem} o DLL que se carga dinámicamente en tiempo de ejecución, de
modo que la extensión queda desacoplada y puede desarrollarse y desplegarse de forma
independiente \parencite[244]{martin2017}. La Figura~\ref{fig:p5-solucion} muestra la
estructura propuesta.

\begin{figure}[!t]
  \centering
  \resizebox{\ifdim\width>\linewidth\linewidth\else\width\fi}{!}{%
  \begin{tikzpicture}[
      caja/.style={draw=azul, fill=azul!6, rounded corners=3pt, thick, align=center,
                   font=\small},
      abs/.style={draw=azul, fill=white, rounded corners=2pt, font=\scriptsize\ttfamily,
                  inner sep=3pt, align=center, minimum width=2.2cm},
      ext/.style={draw=gris, dashed, fill=gris!8, rounded corners=2pt, font=\scriptsize,
                  inner sep=4pt, align=center, minimum width=4.2cm},
      ada/.style={draw=naranja, fill=naranja!10, rounded corners=2pt, font=\scriptsize,
                  inner sep=4pt, align=center, minimum width=4.2cm},
      impl/.style={-{Triangle[open,length=2.4mm,width=2.2mm]}, thick, azul},
      usa/.style={-{Stealth[length=2.2mm]}, thick, gris},
      crea/.style={-{Stealth[length=2.2mm]}, thick, dashed, gris}
    ]
    % Capa externa: interfaz de usuario y Main
    \node[caja, minimum width=14cm, minimum height=1.1cm] (main) at (0,4.35)
      {\textbf{TaxiUI y Main}: configuración y factorías que crean los plugins\\
       \footnotesize(se ajustan al añadir \texttt{Kittens})};
    % Núcleo
    \node[caja, fill=azul!18, minimum width=10.6cm, minimum height=2.3cm] (core) at (0,1.8) {};
    \node[font=\footnotesize\bfseries, azul] at (0,2.6)
      {Núcleo de despacho: políticas estables, cerrado a la modificación};
    \node[abs] at (-3.9,1.55) {«abstracta»\\Buscador};
    \node[abs] at (-1.3,1.55) {«abstracta»\\Selector};
    \node[abs] at (1.3,1.55)  {«abstracta»\\Despachador};
    \node[abs] at (3.9,1.55)  {«interfaz»\\IReglaSeleccion};
    % Plugins
    \node[caja, minimum width=5.4cm, minimum height=1.5cm] (rides) at (-3.7,-1.2)
      {\textbf{Rides} (plugin)\\[1pt]\scriptsize búsqueda y selección de viajes de pasajeros};
    \node[caja, minimum width=5.4cm, minimum height=1.5cm] (kit) at (3.7,-1.2)
      {\textbf{Kittens} (plugin)\\[1pt]\scriptsize selector de entregas · reglas de alergias\\
       \scriptsize registro de alergias · historial de entregas};
    % Adaptadores y sistemas externos
    \node[ada] (a1) at (-3.7,-3.35) {Adaptadores de proveedores de taxis\\(uno por proveedor)};
    \node[ada] (a2) at (3.7,-3.35)  {Adaptadores de Kittens\\(puntos de recogida, proveedores participantes)};
    \node[ext] (x1) at (-3.7,-5.0)  {Sistemas de los proveedores};
    \node[ext] (x2) at (3.7,-5.0)   {Puntos de recogida y proveedores};
    % Relaciones
    \draw[impl] (rides.north) -- (rides.north |- core.south);
    \draw[impl] (kit.north)   -- (kit.north |- core.south);
    \draw[impl] (a1.west) -- ++(-1.6,0) |- (core.west);
    \draw[impl] (a2.north) -- (kit.south);
    \draw[usa]  (a1) -- (x1);
    \draw[usa]  (a2) -- (x2);
    \draw[usa]  (main.south) -- (core.north);
    \draw[crea] ([xshift=-6.6cm]main.south) |- (rides.west);
    \draw[crea] ([xshift=6.6cm]main.south)  |- (kit.east);
    % Leyenda
    \draw[impl] (-6.9,-6.2) -- ++(0.8,0) node[right, font=\scriptsize, black]
      {implementa/extiende (apunta hacia lo estable)};
    \draw[usa]  (0.6,-6.2) -- ++(0.8,0) node[right, font=\scriptsize, black] {usa};
    \draw[crea] (2.9,-6.2) -- ++(0.8,0) node[right, font=\scriptsize, black]
      {crea (factorías en Main)};
  \end{tikzpicture}%
  }
  \caption{Solución propuesta (vista lógica): \texttt{Rides} y \texttt{Kittens} se enchufan
           al núcleo como \emph{plugins}, y los adaptadores aíslan los sistemas externos.
           Entre núcleo y plugins, las dependencias apuntan hacia el núcleo; solo
           \texttt{Main}, el componente más externo, conoce los plugins concretos.
           Elaboración propia a partir de \textcite[244--246]{martin2017} y
           \textcite[21--23]{richards2015}.}
  \label{fig:p5-solucion}
\end{figure}

\paragraph{Por qué respeta el OCP.}
El OCP, formulado por Meyer en 1988 (citado en \cite[70]{martin2017}), establece que un
artefacto de software debe estar abierto a la extensión y cerrado a la modificación. A
nivel arquitectónico se consigue separando la funcionalidad según cómo, por qué y cuándo
cambia, y organizándola en una jerarquía de componentes en la que los de nivel superior
quedan protegidos de los cambios en los de nivel inferior \parencite[74]{martin2017}. En
la propuesta:
\begin{itemize}
  \item el núcleo está \emph{cerrado}: añadir \texttt{Kittens} no modifica ninguna de sus
        clases;
  \item el sistema está \emph{abierto}: la nueva capacidad entra como un componente nuevo
        que extiende las abstracciones del núcleo. Martin lo plantea en términos de
        despliegue: una nueva funcionalidad es un nuevo \emph{jar} que extiende las clases
        abstractas y se añade a la ruta de carga del servicio, lo que se ajusta al OCP
        \parencite[246]{martin2017}.
\end{itemize}

Esta estructura coincide con el patrón \emph{microkernel}, también llamado de
\emph{plugins}: un sistema núcleo al que se añaden módulos \emph{plugin} independientes
que aportan nuevas capacidades de negocio, y que el núcleo localiza mediante un registro
\parencite[21--22]{richards2015}. Martin propone además que la creación y la
configuración se concentren en el componente \texttt{Main}, el más «sucio» de todos
\parencite[232]{martin2017}, que es en sí mismo un \emph{plugin} de la aplicación. Puede
haber varios \texttt{Main}, uno por configuración, por ejemplo, uno por país, jurisdicción
o cliente \parencite[232, 237]{martin2017}. Aplicado a este caso, eso permitiría activar
\texttt{Kittens} solo donde exista el servicio.

\paragraph{Ejemplo: las reglas de alergia como estrategias.}
El requisito de las alergias muestra hasta dónde llega la extensión
(Listado~\ref{lst:p5-reglas}). El núcleo aplica todas las reglas de selección
registradas sin conocerlas, y solo expone identificadores que ya tenía (de vehículo, de
conductor y de cliente). \texttt{Kittens} aporta dos reglas:
\begin{itemize}
  \item excluir a los conductores alérgicos en las entregas;
  \item no asignar a un cliente alérgico un vehículo que haya llevado gatitos en los
        últimos tres días.
\end{itemize}
La segunda regla afecta también a los viajes de pasajeros y, aun así, ni \texttt{Rides}
ni las entidades del núcleo se modifican: los datos de alergias y el historial de
entregas pertenecen al modelo de \texttt{Kittens}, que los consulta por identificador.
Que el núcleo acepte reglas adicionales es la decisión de cierre estratégico que hace
posible la extensión (criterio propio a partir de \cite[106]{martin2017}).

\begin{lstlisting}[style=codigo, language={[Sharp]C}, float=!tb,
  caption={Esbozo en estilo C\# del punto de extensión del núcleo y de las reglas que aporta el plugin Kittens.},
  label={lst:p5-reglas}]
// --- Núcleo de despacho: abstracciones estables (cerradas a la modificación) ---
public interface IReglaSeleccion {
    bool Admite(Vehiculo vehiculo, Pedido pedido);
}

public abstract class Selector {                   // Template Method
    private readonly IEnumerable<IReglaSeleccion> reglas;   // registradas por los plugins
    protected Selector(IEnumerable<IReglaSeleccion> reglas) { this.reglas = reglas; }

    public Vehiculo? Elegir(IEnumerable<Vehiculo> candidatos, Pedido pedido) {
        var admitidos = candidatos.Where(v => reglas.All(r => r.Admite(v, pedido)));
        return Ordenar(admitidos, pedido).FirstOrDefault();   // null si no hay candidatos
    }
    protected abstract IEnumerable<Vehiculo> Ordenar(IEnumerable<Vehiculo> c, Pedido p);
}

// --- Plugin Kittens: modelo y reglas propios; el núcleo no se modifica ---
public class PedidoGatito : Pedido { /* punto de recogida, destino... */ }

public interface IRegistroAlergias {              // datos de Kittens, por identificador
    bool ConductorAlergico(string conductorId);
    bool ClienteAlergico(string clienteId);
}

public class ReglaConductorNoAlergico : IReglaSeleccion {
    private readonly IRegistroAlergias alergias;
    public ReglaConductorNoAlergico(IRegistroAlergias a) { alergias = a; }
    public bool Admite(Vehiculo v, Pedido p) =>
        p is not PedidoGatito || !alergias.ConductorAlergico(v.ConductorId);
}

public class ReglaClienteAlergico : IReglaSeleccion {
    private readonly IRegistroAlergias alergias;
    private readonly IHistorialEntregasGatitos historial;
    public ReglaClienteAlergico(IRegistroAlergias a, IHistorialEntregasGatitos h)
        { alergias = a; historial = h; }
    public bool Admite(Vehiculo v, Pedido p) =>
        !alergias.ClienteAlergico(p.ClienteId)
        || !historial.HuboEntregaEnLosUltimosDias(v.Id, dias: 3);
}
\end{lstlisting}

\paragraph{El papel de los adaptadores.}
El capítulo de Martin sobre el Kitty Problem no menciona los adaptadores. Su papel en esta
extensión se argumenta, como interpretación propia, a partir de la bibliografía:
\begin{itemize}
  \item \textbf{Hacia los sistemas externos.} Un adaptador específico de cada tecnología
        traduce entre el mundo exterior y el puerto de la aplicación, que desconoce lo que
        hay al otro lado \parencite[«Intent»]{cockburn2005}. En la Arquitectura Limpia es
        la capa de adaptadores de interfaz, que convierte los datos entre el formato de
        los casos de uso y las entidades y el de los agentes externos, como otros
        servicios \parencite[205]{martin2017}. La Guía de Arquitectura N-Capas los llama
        Agentes de Servicios externos: aíslan las particularidades de cada servicio, de
        modo que pueda sustituirse sin afectar al sistema
        \parencite[58, 108]{delatorre2010}. \texttt{Kittens} necesita conocer los puntos
        de recogida y qué proveedores participan; esos datos llegan a través de sus
        propios adaptadores, de modo que ni el núcleo ni los adaptadores de proveedores
        existentes cambian, y ni el núcleo ni el plugin conocen formatos o protocolos
        externos.
  \item \textbf{Entre un plugin y el núcleo.} Cuando un plugin trae un contrato propio,
        por ejemplo, porque lo desarrolla un tercero, se crea un adaptador entre ese
        contrato y el contrato estándar del núcleo \parencite[23]{richards2015}. Si la
        logística de gatitos la aportara una empresa asociada, un adaptador la conectaría
        al núcleo sin modificarlo.
\end{itemize}
En ambos casos, los adaptadores permiten que la extensión llegue hasta los sistemas
reales sin romper la Regla de Dependencia: el detalle depende de la abstracción, nunca al
revés.

\section{Valoración crítica}
\label{sec:p5-valoracion}

El patrón de núcleo y plugins aporta lo que el Kitty Problem pedía:
\begin{itemize}
  \item \textbf{agilidad alta}, porque los cambios quedan aislados en plugins débilmente
        acoplados y el núcleo se estabiliza pronto \parencite[25]{richards2015};
  \item \textbf{facilidad de despliegue alta}: según la implementación, los plugins pueden
        añadirse en tiempo de ejecución \parencite[25]{richards2015};
  \item \textbf{testabilidad alta}, porque los plugins pueden probarse de forma aislada
        \parencite[26]{richards2015}.
\end{itemize}
Además, favorece el diseño evolutivo: se parte de un núcleo sólido y se añaden
funcionalidades sin cambios importantes en él \parencite[24--25]{richards2015}.

Tiene, sin embargo, costes y límites:
\begin{enumerate}[label=(\arabic*)]
  \item \textbf{El eje de cambio es una apuesta.} El núcleo solo queda cerrado frente a
        los tipos de cambio más comunes que se esperan o que ya se han vivido
        \parencite[106]{martin2017}. Elegir el «tipo de servicio» como eje se justifica
        por un cambio concreto. La división por pasos sí aislaba otros cambios, como un
        nuevo criterio de selección, que en este diseño debe encajar en el punto de
        extensión del núcleo. Si un requisito futuro afecta a un eje que el núcleo no
        abstrajo, habrá que modificarlo. Además, en la solución de Martin \texttt{TaxiUI}
        también cambia \parencite[244]{martin2017}; en esta propuesta, el ajuste se
        concentra en \texttt{TaxiUI} y \texttt{Main}.
  \item \textbf{Acoplamiento entre plugins.} Los plugins deberían ser independientes entre
        sí y comunicarse lo mínimo \parencite[22]{richards2015}. La regla del cliente
        alérgico hace que \texttt{Kittens} influya en la selección de \texttt{Rides}. El
        punto de extensión del núcleo vuelve explícito ese acoplamiento, pero no lo
        elimina, y las pruebas de la selección deben incluir las reglas registradas.
  \item \textbf{Complejidad de desarrollo.} Exige un diseño cuidadoso y control de los
        contratos: versionado, registro de plugins, granularidad y formas de conexión
        \parencite[26]{richards2015}.
  \item \textbf{Escalabilidad.} Muchas implementaciones del patrón se despliegan como una
        sola unidad, con escalabilidad baja, aunque a veces puede lograrse a nivel de
        plugin \parencite[26]{richards2015}. Si se mantienen servicios distribuidos, hay
        que evitar el error de hacerlos demasiado pequeños
        \parencite[cap.~17]{richardsford2020}.
\end{enumerate}

En conclusión, el Kitty Problem no se debió a usar servicios, sino a trazar las fronteras
sobre los pasos del flujo y no sobre el eje de cambio. Un núcleo estable protegido por la
Regla de Dependencia, capacidades de negocio cohesivas con su propio lenguaje
(\texttt{Rides}, \texttt{Kittens}) que se enchufan como \emph{plugins}, y adaptadores que
aíslan lo externo convierten un cambio que afectaba a todo el sistema en una extensión
localizada. Eso es exactamente lo que persigue el OCP.

\section*{Notas de diseño}

\begin{itemize}
  \item \textbf{Fuente del Kitty Problem.} El enunciado lo sitúa en Richards y Ford, pero
        en la bibliografía consultada aparece en \textcite[cap.~27]{martin2017}. Se
        describe el problema tal como se pide y se aclara el origen en una nota al pie.
  \item \textbf{«Un solo componente», con matices.} Se reconoce explícitamente que la
        interfaz de usuario también cambia, en coherencia con la propia solución de Martin
        y con el cierre estratégico.
  \item \textbf{Regla del cliente alérgico.} Se usa para mostrar un caso que en apariencia
        obligaría a modificar \texttt{Rides} y que, con un punto de extensión en el núcleo
        y los datos propios de \texttt{Kittens}, no lo hace. Es un ejemplo propio, no de
        Martin.
  \item \textbf{Adaptadores.} Como el capítulo 27 no los menciona, su papel se fundamenta
        en \textcite{cockburn2005}, \textcite[205]{martin2017},
        \textcite[58, 108]{delatorre2010} y \textcite[23]{richards2015}, y se presenta
        como interpretación propia.
\end{itemize}

\section*{Notas de investigación}
% ============================================================
% Integrante: Ernesto A. Soler Choong
% Fragmento puro: se incluye desde preguntas/p5-kitty-problem/solucion.tex
% mediante % ============================================================
% Integrante: Ernesto A. Soler Choong
% Fragmento puro: se incluye desde preguntas/p5-kitty-problem/solucion.tex
% mediante % ============================================================
% Integrante: Ernesto A. Soler Choong
% Fragmento puro: se incluye desde preguntas/p5-kitty-problem/solucion.tex
% mediante \input{investigacion/p5-kitty-ocp/notas}.
% Título: se omite «en Richards & Ford» del título oficial porque el Kitty
% Fuente de trabajo: notas.md (misma carpeta).
% Citas: páginas impresas; Richards & Ford (2020) no tiene paginación y se
% cita por capítulo/sección. Todas las claves están en core.bib.
% ============================================================

% Grupo local: \emergencystretch evita líneas desbordadas sin afectar a otros capítulos.
\begingroup\emergencystretch=3em \hbadness=2000
\section*{Notas de investigación: Kitty Problem, OCP y Plugin Architecture}

\section{Síntesis}

El \emph{Kitty Problem} es un ejemplo de \textcite[242--243]{martin2017}, no de Richards y
Ford. Parte de un agregador de taxis, presentado antes en el capítulo sobre el Principio
de Sustitución de Liskov \parencite[80]{martin2017}, construido como microservicios
divididos por función: \texttt{TaxiUI}, \texttt{TaxiFinder}, \texttt{TaxiSelector} y
\texttt{TaxiDispatcher} \parencite[242]{martin2017}. Cuando se pide un servicio de entrega
de gatitos, con restricciones por alergias de conductores y clientes, hay que modificar
\emph{todos} los servicios y coordinar su despliegue \parencite[243]{martin2017}. Martin
usa el caso para desmontar dos supuestos beneficios de los servicios: el desacoplamiento
y el desarrollo y despliegue independientes \parencite[240--241]{martin2017}. Su
diagnóstico es que las descomposiciones funcionales son muy vulnerables a los aspectos
transversales \parencite[243]{martin2017}.

La salida que propone no es crear más servicios, sino aplicar SOLID dentro de ellos
\parencite[244]{martin2017}. La lógica genérica queda en clases base abstractas; lo
específico de los viajes se extrae a un componente \texttt{Rides} y la nueva funcionalidad
va a un componente \texttt{Kittens}. Ambos extienden esas abstracciones mediante
\emph{Template Method} o \emph{Strategy}, respetan la Regla de Dependencia y son creados
por factorías controladas por la interfaz de usuario. Solo cambia \texttt{TaxiUI}; el
resto se extiende añadiendo un componente que se carga en tiempo de ejecución
\parencite[244]{martin2017}, lo que es el Principio Abierto/Cerrado aplicado a
componentes \parencite[246]{martin2017}. La lección arquitectónica: las fronteras no
pasan \emph{entre} los servicios, sino \emph{a través} de ellos
\parencite[246--247]{martin2017}.

Completan el cuadro tres ideas. El CCP es el SRP llevado a componentes, y Martin lo
vincula expresamente con el \emph{Kitty Problem} \parencite[105--106]{martin2017}. La
arquitectura de \emph{plugins} aísla las reglas de negocio de los componentes opcionales
o intercambiables \parencite[170--173]{martin2017}. Y
\textcite[caps.~8 y 17]{richardsford2020} aportan la partición por dominio, el antipatrón
\emph{entity trap} y el bounded context como criterio de granularidad.

\section{Conceptos clave}

\begin{table}[htbp]
  \centering
  \footnotesize
  \renewcommand{\arraystretch}{1.2}
  \begin{tabular}{@{}p{0.21\linewidth}p{0.47\linewidth}p{0.24\linewidth}@{}}
    \toprule
    \textbf{Concepto} & \textbf{Definición} & \textbf{Fuente} \\
    \midrule
    \raggedright \emph{Kitty Problem}
      & \raggedright Un requisito transversal obliga a cambiar y desplegar de forma
        coordinada todos los servicios de un sistema descompuesto por función
      & \raggedright \textcite[242--243]{martin2017} \tabularnewline
    \raggedright Falacias de los servicios
      & \raggedright Siguen acoplados por los datos que comparten; en esa medida, su
        desarrollo y despliegue deben coordinarse
      & \raggedright \textcite[240--241]{martin2017} \tabularnewline
    \raggedright OCP
      & \raggedright Abierto a la extensión y cerrado a la modificación (Meyer, citado por
        Martin); en arquitectura, con SRP y DIP
      & \raggedright \textcite[70, 73--74]{martin2017} \tabularnewline
    \raggedright CCP y cierre estratégico
      & \raggedright Reunir lo que cambia por las mismas razones; cerrar frente a los
        cambios más habituales
      & \raggedright \textcite[105--106]{martin2017} \tabularnewline
    \raggedright Arquitectura de \emph{plugins}
      & \raggedright Las reglas de negocio no dependen de los componentes opcionales; las
        fronteras siguen los ejes de cambio
      & \raggedright \textcite[170--173]{martin2017} \tabularnewline
    \raggedright Regla de Dependencia y adaptadores
      & \raggedright Las dependencias apuntan hacia adentro; los adaptadores convierten
        datos entre el núcleo y lo externo
      & \raggedright \textcite[203, 205]{martin2017} \tabularnewline
    \raggedright Servicios basados en componentes
      & \raggedright Estructura interna SOLID; lo nuevo entra como componentes nuevos y las
        fronteras atraviesan los servicios
      & \raggedright \textcite[245--247]{martin2017} \tabularnewline
    \raggedright Patrón microkernel
      & \raggedright Núcleo mínimo más \emph{plugins} independientes, localizados mediante
        un registro; adaptador para contratos propios
      & \raggedright \textcite[21--23]{richards2015} \tabularnewline
    \raggedright Partición por dominio y \emph{entity trap}
      & \raggedright Con partición técnica, un flujo se reparte entre capas; un componente
        por entidad confunde datos con flujos
      & \raggedright \textcite[cap.~8]{richardsford2020} \tabularnewline
    \raggedright Bounded context y granularidad
      & \raggedright Cada servicio modela un dominio o flujo; fronteras por propósito,
        transacciones y coreografía
      & \raggedright \textcite[caps.~7 y 17]{richardsford2020} \tabularnewline
    \bottomrule
  \end{tabular}
  \caption{Conceptos clave de la investigación INV-P5.}
  \label{tab:inv-p5-conceptos}
\end{table}

\section{Puntos de conexión con las preguntas}

\begin{itemize}
  \item \textbf{P5, por qué falló dividir por verbos o entidades:} la vulnerabilidad de la
        descomposición funcional, las falacias de los servicios, la partición técnica y el
        \emph{entity trap}, y la violación del CCP.
  \item \textbf{P5, diseño con DDD y Onion/Clean:} fronteras sobre los ejes de cambio,
        núcleo protegido por la Regla de Dependencia y bounded context como criterio de
        granularidad.
  \item \textbf{P5, \texttt{Kittens} y \texttt{Rides} como \emph{plugins}:} la solución
        de Martin, el OCP aplicado a componentes, el patrón microkernel y los adaptadores.
  \item \textbf{P3 y P4:} el bounded context y la granularidad apoyan la división en
        contextos de P3; las clases base abstractas y las factorías ilustran el DIP de P4.
\end{itemize}


.
% Título: se omite «en Richards & Ford» del título oficial porque el Kitty
% Fuente de trabajo: notas.md (misma carpeta).
% Citas: páginas impresas; Richards & Ford (2020) no tiene paginación y se
% cita por capítulo/sección. Todas las claves están en core.bib.
% ============================================================

% Grupo local: \emergencystretch evita líneas desbordadas sin afectar a otros capítulos.
\begingroup\emergencystretch=3em \hbadness=2000
\section*{Notas de investigación: Kitty Problem, OCP y Plugin Architecture}

\section{Síntesis}

El \emph{Kitty Problem} es un ejemplo de \textcite[242--243]{martin2017}, no de Richards y
Ford. Parte de un agregador de taxis, presentado antes en el capítulo sobre el Principio
de Sustitución de Liskov \parencite[80]{martin2017}, construido como microservicios
divididos por función: \texttt{TaxiUI}, \texttt{TaxiFinder}, \texttt{TaxiSelector} y
\texttt{TaxiDispatcher} \parencite[242]{martin2017}. Cuando se pide un servicio de entrega
de gatitos, con restricciones por alergias de conductores y clientes, hay que modificar
\emph{todos} los servicios y coordinar su despliegue \parencite[243]{martin2017}. Martin
usa el caso para desmontar dos supuestos beneficios de los servicios: el desacoplamiento
y el desarrollo y despliegue independientes \parencite[240--241]{martin2017}. Su
diagnóstico es que las descomposiciones funcionales son muy vulnerables a los aspectos
transversales \parencite[243]{martin2017}.

La salida que propone no es crear más servicios, sino aplicar SOLID dentro de ellos
\parencite[244]{martin2017}. La lógica genérica queda en clases base abstractas; lo
específico de los viajes se extrae a un componente \texttt{Rides} y la nueva funcionalidad
va a un componente \texttt{Kittens}. Ambos extienden esas abstracciones mediante
\emph{Template Method} o \emph{Strategy}, respetan la Regla de Dependencia y son creados
por factorías controladas por la interfaz de usuario. Solo cambia \texttt{TaxiUI}; el
resto se extiende añadiendo un componente que se carga en tiempo de ejecución
\parencite[244]{martin2017}, lo que es el Principio Abierto/Cerrado aplicado a
componentes \parencite[246]{martin2017}. La lección arquitectónica: las fronteras no
pasan \emph{entre} los servicios, sino \emph{a través} de ellos
\parencite[246--247]{martin2017}.

Completan el cuadro tres ideas. El CCP es el SRP llevado a componentes, y Martin lo
vincula expresamente con el \emph{Kitty Problem} \parencite[105--106]{martin2017}. La
arquitectura de \emph{plugins} aísla las reglas de negocio de los componentes opcionales
o intercambiables \parencite[170--173]{martin2017}. Y
\textcite[caps.~8 y 17]{richardsford2020} aportan la partición por dominio, el antipatrón
\emph{entity trap} y el bounded context como criterio de granularidad.

\section{Conceptos clave}

\begin{table}[htbp]
  \centering
  \footnotesize
  \renewcommand{\arraystretch}{1.2}
  \begin{tabular}{@{}p{0.21\linewidth}p{0.47\linewidth}p{0.24\linewidth}@{}}
    \toprule
    \textbf{Concepto} & \textbf{Definición} & \textbf{Fuente} \\
    \midrule
    \raggedright \emph{Kitty Problem}
      & \raggedright Un requisito transversal obliga a cambiar y desplegar de forma
        coordinada todos los servicios de un sistema descompuesto por función
      & \raggedright \textcite[242--243]{martin2017} \tabularnewline
    \raggedright Falacias de los servicios
      & \raggedright Siguen acoplados por los datos que comparten; en esa medida, su
        desarrollo y despliegue deben coordinarse
      & \raggedright \textcite[240--241]{martin2017} \tabularnewline
    \raggedright OCP
      & \raggedright Abierto a la extensión y cerrado a la modificación (Meyer, citado por
        Martin); en arquitectura, con SRP y DIP
      & \raggedright \textcite[70, 73--74]{martin2017} \tabularnewline
    \raggedright CCP y cierre estratégico
      & \raggedright Reunir lo que cambia por las mismas razones; cerrar frente a los
        cambios más habituales
      & \raggedright \textcite[105--106]{martin2017} \tabularnewline
    \raggedright Arquitectura de \emph{plugins}
      & \raggedright Las reglas de negocio no dependen de los componentes opcionales; las
        fronteras siguen los ejes de cambio
      & \raggedright \textcite[170--173]{martin2017} \tabularnewline
    \raggedright Regla de Dependencia y adaptadores
      & \raggedright Las dependencias apuntan hacia adentro; los adaptadores convierten
        datos entre el núcleo y lo externo
      & \raggedright \textcite[203, 205]{martin2017} \tabularnewline
    \raggedright Servicios basados en componentes
      & \raggedright Estructura interna SOLID; lo nuevo entra como componentes nuevos y las
        fronteras atraviesan los servicios
      & \raggedright \textcite[245--247]{martin2017} \tabularnewline
    \raggedright Patrón microkernel
      & \raggedright Núcleo mínimo más \emph{plugins} independientes, localizados mediante
        un registro; adaptador para contratos propios
      & \raggedright \textcite[21--23]{richards2015} \tabularnewline
    \raggedright Partición por dominio y \emph{entity trap}
      & \raggedright Con partición técnica, un flujo se reparte entre capas; un componente
        por entidad confunde datos con flujos
      & \raggedright \textcite[cap.~8]{richardsford2020} \tabularnewline
    \raggedright Bounded context y granularidad
      & \raggedright Cada servicio modela un dominio o flujo; fronteras por propósito,
        transacciones y coreografía
      & \raggedright \textcite[caps.~7 y 17]{richardsford2020} \tabularnewline
    \bottomrule
  \end{tabular}
  \caption{Conceptos clave de la investigación INV-P5.}
  \label{tab:inv-p5-conceptos}
\end{table}

\section{Puntos de conexión con las preguntas}

\begin{itemize}
  \item \textbf{P5, por qué falló dividir por verbos o entidades:} la vulnerabilidad de la
        descomposición funcional, las falacias de los servicios, la partición técnica y el
        \emph{entity trap}, y la violación del CCP.
  \item \textbf{P5, diseño con DDD y Onion/Clean:} fronteras sobre los ejes de cambio,
        núcleo protegido por la Regla de Dependencia y bounded context como criterio de
        granularidad.
  \item \textbf{P5, \texttt{Kittens} y \texttt{Rides} como \emph{plugins}:} la solución
        de Martin, el OCP aplicado a componentes, el patrón microkernel y los adaptadores.
  \item \textbf{P3 y P4:} el bounded context y la granularidad apoyan la división en
        contextos de P3; las clases base abstractas y las factorías ilustran el DIP de P4.
\end{itemize}


.
% Título: se omite «en Richards & Ford» del título oficial porque el Kitty
% Fuente de trabajo: notas.md (misma carpeta).
% Citas: páginas impresas; Richards & Ford (2020) no tiene paginación y se
% cita por capítulo/sección. Todas las claves están en core.bib.
% ============================================================

% Grupo local: \emergencystretch evita líneas desbordadas sin afectar a otros capítulos.
\begingroup\emergencystretch=3em \hbadness=2000
\section*{Notas de investigación: Kitty Problem, OCP y Plugin Architecture}

\section{Síntesis}

El \emph{Kitty Problem} es un ejemplo de \textcite[242--243]{martin2017}, no de Richards y
Ford. Parte de un agregador de taxis, presentado antes en el capítulo sobre el Principio
de Sustitución de Liskov \parencite[80]{martin2017}, construido como microservicios
divididos por función: \texttt{TaxiUI}, \texttt{TaxiFinder}, \texttt{TaxiSelector} y
\texttt{TaxiDispatcher} \parencite[242]{martin2017}. Cuando se pide un servicio de entrega
de gatitos, con restricciones por alergias de conductores y clientes, hay que modificar
\emph{todos} los servicios y coordinar su despliegue \parencite[243]{martin2017}. Martin
usa el caso para desmontar dos supuestos beneficios de los servicios: el desacoplamiento
y el desarrollo y despliegue independientes \parencite[240--241]{martin2017}. Su
diagnóstico es que las descomposiciones funcionales son muy vulnerables a los aspectos
transversales \parencite[243]{martin2017}.

La salida que propone no es crear más servicios, sino aplicar SOLID dentro de ellos
\parencite[244]{martin2017}. La lógica genérica queda en clases base abstractas; lo
específico de los viajes se extrae a un componente \texttt{Rides} y la nueva funcionalidad
va a un componente \texttt{Kittens}. Ambos extienden esas abstracciones mediante
\emph{Template Method} o \emph{Strategy}, respetan la Regla de Dependencia y son creados
por factorías controladas por la interfaz de usuario. Solo cambia \texttt{TaxiUI}; el
resto se extiende añadiendo un componente que se carga en tiempo de ejecución
\parencite[244]{martin2017}, lo que es el Principio Abierto/Cerrado aplicado a
componentes \parencite[246]{martin2017}. La lección arquitectónica: las fronteras no
pasan \emph{entre} los servicios, sino \emph{a través} de ellos
\parencite[246--247]{martin2017}.

Completan el cuadro tres ideas. El CCP es el SRP llevado a componentes, y Martin lo
vincula expresamente con el \emph{Kitty Problem} \parencite[105--106]{martin2017}. La
arquitectura de \emph{plugins} aísla las reglas de negocio de los componentes opcionales
o intercambiables \parencite[170--173]{martin2017}. Y
\textcite[caps.~8 y 17]{richardsford2020} aportan la partición por dominio, el antipatrón
\emph{entity trap} y el bounded context como criterio de granularidad.

\section{Conceptos clave}

\begin{table}[htbp]
  \centering
  \footnotesize
  \renewcommand{\arraystretch}{1.2}
  \begin{tabular}{@{}p{0.21\linewidth}p{0.47\linewidth}p{0.24\linewidth}@{}}
    \toprule
    \textbf{Concepto} & \textbf{Definición} & \textbf{Fuente} \\
    \midrule
    \raggedright \emph{Kitty Problem}
      & \raggedright Un requisito transversal obliga a cambiar y desplegar de forma
        coordinada todos los servicios de un sistema descompuesto por función
      & \raggedright \textcite[242--243]{martin2017} \tabularnewline
    \raggedright Falacias de los servicios
      & \raggedright Siguen acoplados por los datos que comparten; en esa medida, su
        desarrollo y despliegue deben coordinarse
      & \raggedright \textcite[240--241]{martin2017} \tabularnewline
    \raggedright OCP
      & \raggedright Abierto a la extensión y cerrado a la modificación (Meyer, citado por
        Martin); en arquitectura, con SRP y DIP
      & \raggedright \textcite[70, 73--74]{martin2017} \tabularnewline
    \raggedright CCP y cierre estratégico
      & \raggedright Reunir lo que cambia por las mismas razones; cerrar frente a los
        cambios más habituales
      & \raggedright \textcite[105--106]{martin2017} \tabularnewline
    \raggedright Arquitectura de \emph{plugins}
      & \raggedright Las reglas de negocio no dependen de los componentes opcionales; las
        fronteras siguen los ejes de cambio
      & \raggedright \textcite[170--173]{martin2017} \tabularnewline
    \raggedright Regla de Dependencia y adaptadores
      & \raggedright Las dependencias apuntan hacia adentro; los adaptadores convierten
        datos entre el núcleo y lo externo
      & \raggedright \textcite[203, 205]{martin2017} \tabularnewline
    \raggedright Servicios basados en componentes
      & \raggedright Estructura interna SOLID; lo nuevo entra como componentes nuevos y las
        fronteras atraviesan los servicios
      & \raggedright \textcite[245--247]{martin2017} \tabularnewline
    \raggedright Patrón microkernel
      & \raggedright Núcleo mínimo más \emph{plugins} independientes, localizados mediante
        un registro; adaptador para contratos propios
      & \raggedright \textcite[21--23]{richards2015} \tabularnewline
    \raggedright Partición por dominio y \emph{entity trap}
      & \raggedright Con partición técnica, un flujo se reparte entre capas; un componente
        por entidad confunde datos con flujos
      & \raggedright \textcite[cap.~8]{richardsford2020} \tabularnewline
    \raggedright Bounded context y granularidad
      & \raggedright Cada servicio modela un dominio o flujo; fronteras por propósito,
        transacciones y coreografía
      & \raggedright \textcite[caps.~7 y 17]{richardsford2020} \tabularnewline
    \bottomrule
  \end{tabular}
  \caption{Conceptos clave de la investigación INV-P5.}
  \label{tab:inv-p5-conceptos}
\end{table}

\section{Puntos de conexión con las preguntas}

\begin{itemize}
  \item \textbf{P5, por qué falló dividir por verbos o entidades:} la vulnerabilidad de la
        descomposición funcional, las falacias de los servicios, la partición técnica y el
        \emph{entity trap}, y la violación del CCP.
  \item \textbf{P5, diseño con DDD y Onion/Clean:} fronteras sobre los ejes de cambio,
        núcleo protegido por la Regla de Dependencia y bounded context como criterio de
        granularidad.
  \item \textbf{P5, \texttt{Kittens} y \texttt{Rides} como \emph{plugins}:} la solución
        de Martin, el OCP aplicado a componentes, el patrón microkernel y los adaptadores.
  \item \textbf{P3 y P4:} el bounded context y la granularidad apoyan la división en
        contextos de P3; las clases base abstractas y las factorías ilustran el DIP de P4.
\end{itemize}



